\section{Guide to the analytical proof}\label{sec:analytic-strategy}

This section explains how the three analytical parts fit together.
It may be skipped on a first reading: the definitions and proofs begin
afresh in Section~\ref{sec:diffuse}.

\subsection{Finite density and the diffuse defect}

The main quantities measure departure from a planar transition.
For $|v|<1$, put
\[
 q=1-v^2,\qquad U=g^{-1}(v),\qquad D=1-|\nabla U|^2.
\]
Modica's gradient bound gives $D\ge0$ \cite{Modica1985}.
A planar transition has affine $U$ and $D=0$.
For a stable solution the amplitude $q$ gives the exact identity
\begin{equation}\label{eq:intro-defect-test}
 0\le Q_v(q\eta)
   =\int_{\R^n}q^2|\nabla\eta|^2\dd x
     -\int_{\R^n}q^3D\eta^2\dd x
 \qquad(\eta\in C_c^1(\R^n)).
\end{equation}
The inverse-profile formulation and this test also occur in
\citet[Lemma~2.1]{LLWWW2026}; we derive them in the present dimension.
The second term detects defect, including at zeros where a smooth
level-set description is unavailable.

At a regular point of $\Sigma=\{v=0\}$, let
$\nu=\nabla U/|\nabla U|$ be the normal pointing toward increasing $v$,
let $A=d\nu|_{T\Sigma}$ be the second fundamental form, and set
$H=\operatorname{tr}A$.
Small defect gives a local description by one nearly planar transition.
For its normal averages the appropriate normalization is
\[
 G(s)=1-g(s)^2,\qquad
 \sigma=\int_\R G(s)^2\dd s=\frac{4\sqrt2}{3}=2s_0.
\]
For a fixed, sufficiently large normal length $P$, define
\begin{equation}\label{eq:intro-normal-average}
 I_P(z)=\frac1\sigma\int_{-P}^{P}
 q(z+s\nu(z))^3D(z+s\nu(z))\dd s.
\end{equation}
This is an average along a normal segment, not the pointwise defect
$D(z)$.
We prove that $I_P$ controls both curvature and mean curvature,
with a sharp coefficient in the latter bound.
The limiting equation obtained by normalizing $D$ by its value
$D(z_j)\to0$ at a zero has a nonnegative source even in the normal
coordinate.
Separating this source from the positive homogeneous solutions
gives the strict coefficient needed when the interface has
dimension six.

For the later unbounded-density argument in dimension seven, we will also need a surface inequality. Lifting a surface test into the ambient equation yields
\begin{equation}\label{eq:intro-surface-stability}
 (1-\varepsilon)\int_\Sigma
       \frac{|A|^2+I_P}{2}\psi^2\dd S
 \le \int_\Sigma|\nabla_\Sigma\psi|^2\dd S.
\end{equation}
Here $P$ is chosen first, sufficiently large depending on
$\varepsilon>0$, and $\psi$ is Lipschitz and compactly supported
where $D<\delta$, with $\delta$ then sufficiently small.
The two favorable terms come from normal tube geometry and profile
defect, respectively.  When interaction dominates curvature, this
inequality retains only about half of the mass $I_P$.
The final ambient test will retain almost all of that mass for
the specifically chosen radial amplitude.  This improvement is what makes
the eventual comparison in dimension seven decisive.

Returning to finite density, we use the preliminary critical-density and regularity results of Florit-Simon and Serra \cite[Sections~3--5]{FSS2025}.
These results apply through ambient dimension seven; their final
classification theorem is four-dimensional.
They reduce a possible failure of Theorem~\ref{thm:finite} to
an integer-density solution whose zero set consists, away from a
sparse set, of ordered sheets.
The transition-layer estimates of Wang and Wei
\cite{WangWei2019} apply after establishing the required
nonvanishing core gradient and a bound on the full ambient
derivative of its direction.
We keep these hypotheses explicit.

The remaining contradiction compares the nearest gap between
sheets in two ways.  A preliminary concave test controls the
geometric errors.  A sharper gap test then combines the
mean-curvature bound with a radial Green estimate, while diffuse
stability supplies the competing Hardy bound.
Their logarithmic coefficients are incompatible through
dimension seven.  The resulting finite-density theorem makes
every fixed finite density threshold available in the local
regularity theorem used later.

\subsection{Geometry without a global density bound}

The remaining task is to exclude unbounded density for a stable
solution in $\R^7$.
The relevant geometry can first be seen in an ordinary normal
tube around a smooth sheet: its points have the form
$z+t\nu(z)$, where $\nu$ is the unit normal.
Such a tube can cease to describe one isolated transition because
the sheet bends or because a second sheet enters it.
Both obstructions must be measured when the number of sheets
is not known in advance.

For two different zeros $z,y$, with $z$ regular, the ratio
\[
\frac{|y-z|^2}{|\nu(z)\cdot(y-z)|}
\]
is twice the radius of the normal tangent ball through $y$.
When the denominator vanishes, the ratio is interpreted as
$+\infty$ and there is no such finite ball.
As $y$ approaches $z$, it also records curvature.
We use a smoothed angular penalty to define a comparable length
$d(z)$, called the \emph{penalized reach}.
Its purpose is to retain these two geometric controls while
allowing a differential comparison at a second contact point.
The penalty and its precise properties are introduced where
this comparison is proved.

A delayed energy-monotonicity argument finds a bounded-density
window at nearly the reach scale.  The preceding local regularity
theorem then isolates a single sheet there, although the original
solution has no global density bound.
Projection onto the translation mode of $g$ improves the
mean-curvature derivatives.  A tangent ball on one phase side
controls signed mean curvature at a remote contact.
These estimates produce a differential inequality for $d$ and
a change-of-variables formula that counts its positive curvature
contributions with multiplicity one.

We next connect this local geometry to stability.
Tests lifted from a sheet use the derivative of the
one-dimensional profile.  Where several lifts reach the same
ambient point, their squares are added before taking a square
root.  The resulting gradient identity retains a nonnegative
variance term.
Joining these tests to $q$ near zeros with large defect preserves
positive defect mass and controls the cost of the joining region.
Combining the resulting inequality with the
Simons--Schoen--Simon--Yau curvature method
\cite{Simons1968,SchoenSimonYau1975} yields a fourth-moment growth
bound.  Interaction is measured here by $E=e^{-\kappa d}$,
where $\kappa>0$ is fixed by the reach penalty.
The estimate controls weighted integrals of $|A|^4+D^2+E^2$,
together with the weighted count of a separated net of large-defect zeros, by
$K(R)R^{m-4}$, where $m=6$ and
$K(R)=1+\sup_xM_R(x)$.
This gains four powers of $R$ over the corresponding weighted
area bound, even though $K(R)$ may grow.
The mean curvature need not vanish, so its derivative terms
remain part of this calculation.

This passage from ambient stability to the transition set is
related to the free-boundary work of Chan, Fern\'andez-Real,
Figalli and Serra \cite{CFFS2025} and to
\citet[Sections~3 and~5]{LLWWW2026} for the smooth quartic
equation in dimension three.
The latter joins $1-v^2$ to transition-profile derivatives and
uses a surface inequality, Gauss--Bonnet, and logarithmic cutoffs.
Here the six-dimensional interface requires the reach inequality,
the fourth-moment growth estimate, and the final radial comparison.
The two joining operations also have different purposes:
the first changes from ambient to surface tests near large-defect
zeros; the final one interpolates amplitudes between neighboring
sheets.

In that final comparison, the same radial reach weight is used
in the geometric inequality and the ambient stability test.
The growth estimate controls the regions where the regular-sheet
description fails.
On the remaining regions, the surface amplitude is extended normally
and interpolated between adjacent sheets in fixed-width bands
inside their logarithmic gaps.
The cost is proportional to the squared difference of the two
amplitudes, so it vanishes when they agree.
Assign each portion of a switching band to the base with the larger
uncut amplitude.  On that base's nearest gap, the reach comparison
makes the amplitude ratio tend to one, so the normalized loss vanishes.
Of the two possible sides at the assigned base, only the other side
can therefore incur the fixed positive loss.
The resulting stability estimate retains almost nine tenths of
$\int_\Sigma I_P\psi^2\dd S$, while the coefficient of
$\int_\Sigma|\nabla_\Sigma\psi|^2\dd S$ stays close to one.
For the same radial amplitude $\psi$, this is stronger in the
interaction term than \eqref{eq:intro-surface-stability} and
provides the additional control of $I_P$ needed in the reach comparison.
Expanding the shared radial gradient and comparing its coefficients
with the reach inequality leaves a strict margin in interface dimension
six.  A normal-tube packing argument then converts the
estimate into an area contradiction, excluding unbounded density
and completing Theorem~\ref{thm:stable}.

